\scanpage{73}
On va prouver la proposition auxiliaire

\uline{Prop 3} On a, pour $u\in L_0(E,F)$, et $0\leq m\leq n$

\[
\begin{array}{@{}l@{\qquad}l@{}}
\textcircled{19}\ \deleted{\textcircled{18}} &
\displaystyle \sigma_n^m(u)=\operatorname*{Sup}_{\left\{\substack{v_n\in L(F,E)\\\lVert v_n\rVert\leq1\\\operatorname{rang}v_n\leq n}\right.}|\alpha_m(v_nu)|=\operatorname*{Sup}_{v_n}\{\alpha_{\unsure{n|m}}(|v_nu|)
\end{array}
\]

\deleted{\unsure{En effet}}.

En effet, prouvons, \deleted{que $|\alpha_m(v_nu)|=|\operatorname{Tr}\overset{m}{\Lambda}v_nu|\leq$} \deleted{Il suffit} pour prouver l'inégalité \deleted{$|\operatorname{Tr}(\overset{m}{\Lambda}v_n)(\overset{m}{\Lambda}u)|\leq\lVert\overset{m}{\Lambda}\unsure{v_n}\rVert\,\lVert\overset{m}{\Lambda}u\rVert_1\leq\overset{m}{\Lambda}u$} $|\alpha_m(v_nu)|\leq\underline{\sigma_n^m}(u)$, comme on a $\rho_i(v_nu)\leq\rho_i(u)$ pour tout $i$ (\unsure{---}), il suffit d'établir $|\alpha_m(v_nu)|\leq\sigma_n^m(v_nu)$, donc le

\uline{lemme} Si $u\in L(E,E)$ est de rang $\leq n$, on a $|\alpha_m(u)|\leq\{\alpha_m(|u|)\}\leq\sigma_n^m(u)$. Or on a, si $w$ est la restriction de $u$ à $u(E)$, \deleted{et $\alpha$ l'application identique $u(E)\longrightarrow E$;}

\[
|\alpha_m(\deleted{\unsure{uv}}u)|=|\alpha_m(w)|\leq\deleted{\unsure{\alpha_m(|v|)}}\sum_{1\leq i_1\mathord{-}\leq i_m\leq n}\rho_{i_1}(w)\mathord{-}\rho_{i_m}(w),
\]

et comme \deleted{\unsure{que}} $w=\alpha^*u\alpha$ ($\alpha$ appl. identique $u(E)\longrightarrow E$) d'où $\rho_i(w)\leq\lVert\alpha^*\rVert\lVert\alpha\rVert\rho_i(u)\leq\rho_i(u)$, la formule voulue.

\margin{\unsure{1|$\uparrow$}}

\deleted{\unsure{Montrons}} Montrons que la borne en (\unsure{19|18}) est \uuline{atteinte} : on peut supposer $E$ de dimension $\geq n$, \deleted{\unsure{alors}} et si \deleted{\unsure{l'on est} $\sum_{\unsure{i=1|i}}^{\infty}$ \unsure{des} \unsure{$\rho_i e_i\otimes f_i$}}, $u=\sum\rho_i\bar e_i\otimes f_{\unsure{2|i}}$, on prend $v_n=\sum_1^n\bar f_i\otimes e_i$, d'où $v_nu=\sum_1^n\rho_i\bar e_i\otimes e_i$, et $\alpha_m(v_nu)=S_n^m(u)$ (remarque : si $F=E$, et $u$ hermitien, alors on peut prendre $v$ projecteur hermitien de rang $n$ pour $E$ de dimension $\leq n$).
Prouvons, maintenant (16) : \deleted{\unsure{Or} (17) nous ramène à prouver que} \inserted{(18) nous ramène à prouver} que $|\alpha_m(w(u+v))|\deleted{\unsure{\leq\sum_{p=0}^{n}\alpha_p(wu)}}$ est majoré par le second m\textsuperscript{b} pour \deleted{\unsure{tout}} $w\in L(F,E)$ telle que $\lVert w\rVert\leq1$, $\operatorname{rang}w\leq n$, \deleted{\unsure{or il est}} \deleted{\unsure{on}} on a ((14))

\[
|\alpha_{\unsure{n|m}}(w(u+v))|\leq\alpha_{\unsure{n|m}}(|wu+wv|)\leq\sum_{p=0}^{m}\alpha_p(|wu|)\alpha_{m-p}(|wv|)
\]

\deleted{d'autre part on a}

\[
\deleted{\alpha_p(|wu|)=\lVert\overset{p}{\Lambda}wu\rVert_1=\lVert(\overset{p}{\Lambda}w)(\overset{p}{\Lambda}u)\rVert_1\leq\lVert\overset{\unsure{n|p}}{\Lambda}w\rVert\,\lVert\overset{p}{\Lambda}u\rVert_1\leq\lVert\overset{\unsure{n|p}}{\Lambda}u\rVert_1=\unsure{\alpha_p(|u|)}}
\]

\deleted{or l'opération \unsure{de la}} Or (\unsure{19|18}) montre que $\alpha_p(|wu|)\leq\sigma_n^p(u)$ et $\alpha_{m-p}(|wv|)\leq\sigma_n^{m-p}(v)$, d'où la formule (18),
\margin{74}
